\chapter{Revisão do Plano de Ação}
\label{cap:plano-acao}

Este capítulo compara o plano do Projeto Final I com o que foi realizado, discute as dificuldades e apresenta os resultados preliminares, que sustentam a viabilidade da proposta. A proposta previa a avaliação no CloudSim Plus, mas o protótipo usou um simulador próprio, em Python, porque a decisão estudada se resume a escolher $n_t$ a partir da demanda agregada e o simulador pôde ser coberto por testes; ele continua no Projeto Final II, com o CloudSim Plus como validação opcional. Ficam para o Projeto Final II os objetivos \ref{oe:memoria}, \ref{oe:robustez} e \ref{oe:replicacao} e a conclusão de \ref{oe:custo} a \ref{oe:sensibilidade} (Tabela~\ref{tab:cronograma-pf2}).

\section{Atividades realizadas}
\label{sec:atividades-realizadas}
\label{sec:plano-previsto}

A meta do Projeto Final I foi um protótipo de ponta a ponta, do rastro bruto da Alibaba ao custo e às violações de cada política, como evidência de viabilidade. A construção foi dividida em fases curtas, cada uma com uma especificação e um critério de aceite, todas concluídas antes do prazo de 31/10/2026; as datas, tiradas do histórico de versões do repositório \cite{rosas2026repo}, estão na Tabela~\ref{tab:cronograma-pf1}.

O rastro da Alibaba foi agregado em uma série de demanda de processador com intervalos de cinco minutos, que cobre os oito dias do rastro. Como nem todas as máquinas enviam medições em todos os intervalos, a demanda é a utilização média das máquinas que reportaram, estendida ao total de máquinas, o que evita subestimá-la nos intervalos incompletos. As lacunas curtas foram interpoladas, e a interrupção de cerca de cinco horas na coleta, dentro do treino, foi preenchida com o mesmo horário do dia anterior. A Figura~\ref{fig:serie} mostra a série, a divisão cronológica, o trecho preenchido, em tracejado, e a cobertura, com o limite abaixo do qual o intervalo é lacuna. A validação e o teste têm cobertura quase completa, e o trecho menos confiável, no início do rastro, fica no treino.

\begin{figure}[htbp]
\centering
\caption{Demanda de processador e cobertura do rastro da Alibaba}
\includegraphics[width=0.85\textwidth]{figures/series_full.png}
\label{fig:serie}
\fonte{Elaborado pelo autor.}
\end{figure}

Os seis previsores foram implementados sob a interface comum da Seção~\ref{sec:solucao} e avaliados com origem móvel (\textit{walk-forward}); testes automatizados verificam que alterar os dados posteriores à origem não muda nenhuma previsão. O simulador e as políticas têm testes de casos calculados à mão e verificações de sanidade, como a de que o oráculo não viola nenhum intervalo.

\section{Dificuldades encontradas}
\label{sec:dificuldades}

\textbf{Holt-Winters.} O statsmodels não aplica os parâmetros estimados a dados novos, e reajustar o modelo em cada origem seria lento demais. A recursão do método foi implementada no próprio previsor, com os parâmetros do treino, e um teste confirma que ela reproduz a biblioteca até a precisão de máquina.

\textbf{Ponto flutuante.} Três problemas tiveram origem no ponto flutuante: o oráculo aparecia com violações inexistentes, porque $k \cdot \lceil d_t / k \rceil$ pode ficar ligeiramente abaixo de $d_t$; o arredondamento para cima saltava para o inteiro seguinte quando o valor exato era inteiro; e a soma paralela do DuckDB variava na última casa decimal entre execuções. Comparações e arredondamentos passaram a usar tolerância relativa de $10^{-9}$, e a série, seis casas decimais.

\textbf{Tamanho do teste.} Com 462 intervalos, um intervalo com violação vale 0,22 ponto percentual de $V$, e a diferença entre cumprir ou não $\varepsilon = 1\%$ no teste é de um único intervalo. A primeira versão descartava o início do rastro, de cobertura menor; usar o rastro completo aumentou o teste em 30\%, mas não mudou a ordem de grandeza do problema, o que motiva a avaliação com várias janelas de teste (\ref{oe:robustez}).

\section{Resultados preliminares}
\label{sec:resultados-preliminares}

Os resultados, com os parâmetros da Tabela~\ref{tab:parametros}, são preliminares: com uma única janela de teste e uma única semente da GRU, diferenças de poucos por cento entre previsores não permitem afirmar que um é melhor que outro.

No erro de previsão (Tabela~\ref{tab:erro-previsao}, em máquinas equivalentes), o ARIMA e a GRU são os melhores, e o sazonal ingênuo, que repete o valor de um dia antes, é o pior por larga margem; o erro nos demais horizontes está no Apêndice~\ref{ap:erro-horizonte}.

\begin{table}[htbp]
\centering
\caption{Erro das previsões no teste, com $h = 1$}
\label{tab:erro-previsao}
\footnotesize
\begin{tabular}{lrrrr}
\toprule
\textbf{Previsor} & \textbf{RMSE} & \textbf{MAE} & \textbf{Viés} & \textbf{Taxa de subestimação} \\
\midrule
ARIMA & 157,44 & 120,25 & $+1{,}75$ & 42,9\% \\
GRU & 162,26 & 125,56 & $+17{,}59$ & 39,4\% \\
Ingênuo & 168,05 & 126,83 & $+0{,}56$ & 45,0\% \\
Holt-Winters & 173,68 & 133,44 & $-1{,}13$ & 51,5\% \\
Média móvel & 187,65 & 142,68 & $+3{,}01$ & 46,5\% \\
Sazonal ingênuo & 360,90 & 291,42 & $-166{,}41$ & 66,9\% \\
\bottomrule
\end{tabular}
\fonte{Elaborado pelo autor.}
\end{table}

A Figura~\ref{fig:custo} resume o custo de cada política no seu ponto de operação, detalhado na Tabela~\ref{tab:pontos-operacao}. O principal resultado é que, entre as políticas que cumprem o limite no teste, ajustar a capacidade à demanda reduz o custo em 24\% a 30\% com $\varepsilon = 1\%$ e em 32\% a 36\% com $\varepsilon = 5\%$ em relação ao dimensionamento estático, faixa comparável à de \citeonline{calheiros2015arima}. Esse ganho vem da elasticidade, e não da previsão: em $\varepsilon = 5\%$, a melhor preditiva, com o ARIMA, fica só 1,7\% abaixo da reativa, e as diferenças entre previsores são de poucos por cento, com a GRU e o ARIMA alternando a primeira posição entre os dois limites. O custo acima do oráculo é quase todo capacidade ociosa, já que a demanda não atendida é desprezível em quase todas as políticas (Tabela~\ref{tab:uoa}).

\begin{figure}[htbp]
\centering
\caption{Custo de cada política no ponto de operação}
\includegraphics[width=\textwidth]{figures/cost_at_epsilon.png}
\label{fig:custo}
\fonte{Elaborado pelo autor.}
\end{figure}

\begin{table}[htbp]
\centering
\caption{Pontos de operação no teste}
\label{tab:pontos-operacao}
\footnotesize
\setlength{\tabcolsep}{3pt}
\begin{tabular}{l rrr rr @{\hspace{6pt}} rrr rr}
\toprule
 & \multicolumn{5}{c}{$\varepsilon = 1\%$} & \multicolumn{5}{c}{$\varepsilon = 5\%$} \\
\cmidrule(lr){2-6} \cmidrule(lr){7-11}
 & \multicolumn{3}{c}{$L = 1$} & \multicolumn{2}{c}{$L = 2$} & \multicolumn{3}{c}{$L = 1$} & \multicolumn{2}{c}{$L = 2$} \\
\cmidrule(lr){2-4} \cmidrule(lr){5-6} \cmidrule(lr){7-9} \cmidrule(lr){10-11}
\textbf{Política} & \textbf{Par.} & $C_{\mathrm{norm}}$ & $V$ (\%) & $C_{\mathrm{norm}}$ & $V$ (\%) & \textbf{Par.} & $C_{\mathrm{norm}}$ & $V$ (\%) & $C_{\mathrm{norm}}$ & $V$ (\%) \\
\midrule
Oráculo & -- & 1,000 & 0,00 & 1,000 & 0,00 & -- & 1,000 & 0,00 & 1,000 & 0,00 \\
Dimensionamento estático & -- & 1,830 & 0,22 & 1,830 & 0,22 & -- & 1,830 & 0,22 & 1,830 & 0,22 \\
Reativa (HPA) & 0,82 & 1,340 & \textbf{1,08} & 1,376 & 0,65 & 0,92 & 1,197 & 4,98 & 1,223 & 4,98 \\
\textit{Preditiva com:} & & & & & & & & & & \\
\quad ARIMA & 0,32 & 1,315 & 0,87 & 1,385 & 0,65 & 0,18 & 1,176 & 3,90 & 1,238 & 3,25 \\
\quad GRU & 0,28 & 1,287 & 0,87 & 1,356 & 0,43 & 0,19 & 1,199 & 3,25 & 1,226 & 3,68 \\
\quad Holt-Winters & 0,40 & 1,391 & 0,43 & 1,440 & 0,87 & 0,19 & 1,185 & 3,46 & 1,233 & 3,03 \\
\quad Ingênuo & 0,38 & 1,373 & 0,87 & 1,412 & 0,65 & 0,21 & 1,206 & 3,46 & 1,275 & 3,90 \\
\quad Média móvel & 0,36 & 1,354 & 0,43 & 1,434 & 0,43 & 0,24 & 1,238 & 3,90 & 1,238 & 3,90 \\
\quad Sazonal ingênuo & 0,42 & 1,270 & \textbf{13,64} & 1,270 & \textbf{13,64} & 0,23 & 1,103 & \textbf{33,55} & 1,103 & \textbf{33,55} \\
\bottomrule
\end{tabular}

\fonte{Elaborado pelo autor.}
\smallskip
\parbox{\textwidth}{\footnotesize Nota: Par. é a margem $m$ na política preditiva e a utilização-alvo $u^{*}$ na reativa, escolhida na validação (Seção~\ref{sec:solucao}) com $L = 1$; com $L = 2$, a escolha é refeita da mesma forma. Valores de $V$ em negrito ultrapassam $\varepsilon$ no teste.}
\end{table}

\begin{table}[htbp]
\centering
\caption{Demanda não atendida, ociosidade e alterações no teste}
\label{tab:uoa}
\footnotesize
\begin{tabular}{l rrr @{\hspace{12pt}} rrr}
\toprule
 & \multicolumn{3}{c}{$\varepsilon = 1\%$} & \multicolumn{3}{c}{$\varepsilon = 5\%$} \\
\cmidrule(lr){2-4} \cmidrule(lr){5-7}
\textbf{Política} & $U$ (\%) & $O$ (\%) & $A$ & $U$ (\%) & $O$ (\%) & $A$ \\
\midrule
Oráculo & 0,00 & 1,89 & 378 & 0,00 & 1,89 & 378 \\
Dimensionamento estático & 0,01 & 46,41 & 0 & 0,01 & 46,41 & 0 \\
Reativa (HPA) & 0,06 & 26,84 & 116 & 0,39 & 18,37 & 116 \\
\textit{Preditiva com:} & & & & & & \\
\quad ARIMA & 0,04 & 25,40 & 386 & 0,29 & 16,84 & 379 \\
\quad GRU & 0,06 & 23,83 & 345 & 0,24 & 18,37 & 326 \\
\quad Holt-Winters & 0,07 & 29,50 & 409 & 0,31 & 17,49 & 411 \\
\quad Ingênuo & 0,02 & 28,57 & 403 & 0,25 & 18,82 & 397 \\
\quad Média móvel & 0,06 & 27,60 & 284 & 0,26 & 20,92 & 282 \\
\quad Sazonal ingênuo & 1,22 & 23,66 & 408 & 4,34 & 14,89 & 398 \\
\bottomrule
\end{tabular}

\fonte{Elaborado pelo autor.}
\smallskip
\parbox{\textwidth}{\footnotesize Nota: valores nos pontos de operação da Tabela~\ref{tab:pontos-operacao}; o oráculo e o dimensionamento estático não dependem de $\varepsilon$.}
\end{table}

A ordenação por erro também não se repete na ordenação por custo. A GRU tem RMSE 3\% maior que o do ARIMA, mas custa menos em $\varepsilon = 1\%$, coerente com seu viés positivo e sua menor taxa de subestimação, que pedem margem menor (0,28 contra 0,32). O Holt-Winters é o previsor mais caro em $\varepsilon = 1\%$, em parte porque viola menos ($V$ de 0,43\% contra 0,87\%, dois intervalos), e um dos mais baratos em $\varepsilon = 5\%$. Já o sazonal ingênuo, de pior erro, aparece como o mais barato só porque ultrapassa $\varepsilon$ no teste: a demanda muda de comportamento entre a validação e o teste, e a margem escolhida na validação deixa de cobrir a subestimação. A reativa em $\varepsilon = 1\%$ também passa do limite, por um único intervalo.

Por fim, a política reativa altera o número de instâncias com menos da metade da frequência das preditivas, vantagem que não entra no custo. Com antecedência $L = 2$, a vantagem da previsão sobre a reativa diminui: entre as preditivas que cumprem o limite, só a GRU fica abaixo dela em $\varepsilon = 1\%$, 1,5\% a menos, e nenhuma em $\varepsilon = 5\%$. Isso contraria a motivação de que o atraso prejudica mais a reativa; hipóteses a testar em \ref{oe:sensibilidade} são o RMSE de 10\% a 25\% maior com $h = 2$ (Tabela~\ref{tab:erro-horizonte-test}) e a janela de estabilização da reativa.
